%======================================================================
\section{Data Frames, Factors \& Tables, and File I/O}
%======================================================================

\sectionframe

\begin{frame}[fragile]{Dataframe}
\begin{itemize}
  \item \textit{A two-dimensional table of vectors, that can contain data of different types}
  \item To create a data frame, use \code{data.frame()}:
\end{itemize}
\begin{lstlisting}[language=R]
x <- 1:5
y <- c("M", "F", "M", "M", "F")
z <- c("one", "two", "three", "four", "five")

my_df <- data.frame(x, y, z)
my_df
\end{lstlisting}
\begin{lstlisting}[style=output]
  x y     z
1 1 M   one
2 2 F   two
3 3 M three
4 4 M  four
5 5 F  five
\end{lstlisting}
\end{frame}

\begin{frame}[fragile]{Dataframe: Attributes}
\begin{lstlisting}[language=R]
my_df
\end{lstlisting}
\begin{lstlisting}[style=output]
  x y     z
1 1 M   one
2 2 F   two
3 3 M three
4 4 M  four
5 5 F  five
\end{lstlisting}
\begin{itemize}
  \item Attributes \textit{(similar to matrix)}:
\end{itemize}
\begin{lstlisting}[language=R]
nrow(my_df)  # number of observations
\end{lstlisting}
\begin{lstlisting}[style=output]
[1] 5
\end{lstlisting}
\begin{lstlisting}[language=R]
ncol(my_df)  # number of variables
\end{lstlisting}
\begin{lstlisting}[style=output]
[1] 3
\end{lstlisting}
\end{frame}

\begin{frame}[fragile]{Dataframe: Attributes (cont.)}
\begin{lstlisting}[language=R]
dim(my_df)  # dimensions of the dataframe
\end{lstlisting}
\begin{lstlisting}[style=output]
[1] 5 3
\end{lstlisting}
\begin{itemize}
  \item Other helpful attribute functions:
  \begin{itemize}
    \item \code{head(df, n)} will look at the first \textit{n} observations of a dataframe
    \item \code{tail(df, n)} will look at the last \textit{n} observations of a dataframe
  \end{itemize}
\end{itemize}
\begin{columns}[T]
\begin{column}{0.48\textwidth}
\begin{lstlisting}[language=R]
head(my_df, 3)
\end{lstlisting}
\begin{lstlisting}[style=output]
  x y     z
1 1 M   one
2 2 F   two
3 3 M three
\end{lstlisting}
\end{column}
\begin{column}{0.48\textwidth}
\begin{lstlisting}[language=R]
tail(my_df, 3)
\end{lstlisting}
\begin{lstlisting}[style=output]
  x y     z
3 3 M three
4 4 M  four
5 5 F  five
\end{lstlisting}
\end{column}
\end{columns}
\end{frame}

\begin{frame}[fragile]{Dataframe: Subsetting}
\begin{lstlisting}[language=R]
my_df
\end{lstlisting}
\begin{lstlisting}[style=output]
  x y     z
1 1 M   one
2 2 F   two
3 3 M three
4 4 M  four
5 5 F  five
\end{lstlisting}
\begin{itemize}
  \item Subsetting a data frame:
  \item First, using a pair of indices, like a matrix:
\end{itemize}
\begin{lstlisting}[language=R]
my_df[1, 2]  # extract the element in the 1st row and 2nd col
\end{lstlisting}
\begin{lstlisting}[style=output]
[1] "M"
\end{lstlisting}
\begin{lstlisting}[language=R]
my_df[, 1]  # extract only the 1st col
\end{lstlisting}
\begin{lstlisting}[style=output]
[1] 1 2 3 4 5
\end{lstlisting}
\end{frame}

\begin{frame}[fragile]{Dataframe: Subsetting (cont.)}
\begin{lstlisting}[language=R]
my_df[3, 2:3]  # 3rd row, col 2 through 3
\end{lstlisting}
\begin{lstlisting}[style=output]
  y     z
3 M three
\end{lstlisting}
\begin{lstlisting}[language=R]
my_df[c(1, nrow(my_df)), ]  # extract the 1st and the last obs
\end{lstlisting}
\begin{lstlisting}[style=output]
  x y    z
1 1 M  one
5 5 F five
\end{lstlisting}
\begin{itemize}
  \item Second, by name:
\end{itemize}
\begin{lstlisting}[language=R]
names(my_df)  # use names() to look at the names
\end{lstlisting}
\begin{lstlisting}[style=output]
[1] "x" "y" "z"
\end{lstlisting}
\end{frame}

\begin{frame}[fragile]{Dataframe: Subsetting (cont.)}
\begin{lstlisting}[language=R]
my_df$x  # access variable x
\end{lstlisting}
\begin{lstlisting}[style=output]
[1] 1 2 3 4 5
\end{lstlisting}
\begin{lstlisting}[language=R]
my_df$x[4]  # access the 4th element in variable x
\end{lstlisting}
\begin{lstlisting}[style=output]
[1] 4
\end{lstlisting}
\begin{lstlisting}[language=R]
mean(my_df$x)  # calculate the mean of variable x
\end{lstlisting}
\begin{lstlisting}[style=output]
[1] 3
\end{lstlisting}
\end{frame}

\begin{frame}[fragile]{Miscellaneous Functions for Text Manipulation}
\begin{itemize}
  \item \code{paste()} joins strings; by default, \code{paste()} will use \mbox{\code{sep = " "}}
\end{itemize}
\begin{lstlisting}[language=R]
paste("a", "b")
\end{lstlisting}
\begin{lstlisting}[style=output]
[1] "a b"
\end{lstlisting}
\begin{lstlisting}[language=R]
paste("a", "b", sep = "")
\end{lstlisting}
\begin{lstlisting}[style=output]
[1] "ab"
\end{lstlisting}
\begin{lstlisting}[language=R]
paste0("a", "b")  # a special version of paste() assuming sep = ""
\end{lstlisting}
\begin{lstlisting}[style=output]
[1] "ab"
\end{lstlisting}
\end{frame}

\begin{frame}[fragile]{Miscellaneous Functions for Text Manipulation (cont.)}
\begin{itemize}
  \item Assigning two vectors:
\end{itemize}
\begin{lstlisting}[language=R]
c1 <- c("a", "b")
c2 <- c("c", "d")
\end{lstlisting}
\begin{itemize}
  \item Difference between \code{sep} and \code{collapse}:
  \item \code{paste0(c1, c2)} vs \code{paste0(c1, c2, collapse = "+")}:
\end{itemize}
\begin{lstlisting}[language=R]
paste0(c1, c2)
\end{lstlisting}
\begin{lstlisting}[style=output]
[1] "ac" "bd"
\end{lstlisting}
\begin{lstlisting}[language=R]
paste0(c1, c2, collapse = "+")
\end{lstlisting}
\begin{lstlisting}[style=output]
[1] "ac+bd"
\end{lstlisting}
\end{frame}

\begin{frame}[fragile]{Miscellaneous Functions for Text Manipulation (cont.)}
\begin{itemize}
  \item Using \code{sep} and \code{collapse} together:
\end{itemize}
\begin{lstlisting}[language=R]
paste(c1, c2, sep = "_", collapse = "+")
\end{lstlisting}
\begin{lstlisting}[style=output]
[1] "a_c+b_d"
\end{lstlisting}
\begin{itemize}
  \item Concatenation between variable and text strings:
\end{itemize}
\begin{lstlisting}[language=R]
ans <- 3
paste0("The value of variable ans is ", ans)
\end{lstlisting}
\begin{lstlisting}[style=output]
[1] "The value of variable ans is 3"
\end{lstlisting}
\end{frame}

\begin{frame}[fragile]{Factors}
\begin{itemize}
  \item In R, we usually use factor variable types for categorical variables.
  \item Here, vector \code{f1} is numeric:
\end{itemize}
\begin{lstlisting}[language=R]
f1 <- c(0, 1, 1, 1, 0, 0, 1)
\end{lstlisting}
\begin{lstlisting}[language=R]
class(f1)
\end{lstlisting}
\begin{lstlisting}[style=output]
[1] "numeric"
\end{lstlisting}
\begin{itemize}
  \item We can use the \code{factor()} function to change \code{f1}:
\end{itemize}
\begin{lstlisting}[language=R]
f1 <- factor(f1)
\end{lstlisting}
\begin{columns}[T]
\begin{column}{0.48\textwidth}
\begin{lstlisting}[language=R]
class(f1)
\end{lstlisting}
\begin{lstlisting}[style=output]
[1] "factor"
\end{lstlisting}
\end{column}
\begin{column}{0.48\textwidth}
\begin{lstlisting}[language=R]
f1
\end{lstlisting}
\begin{lstlisting}[style=output]
[1] 0 1 1 1 0 0 1
Levels: 0 1
\end{lstlisting}
\end{column}
\end{columns}
\end{frame}

\begin{frame}[fragile]{Factors (cont.)}
\begin{itemize}
  \item Now that \code{f1} is a factor vector, we can utilize its levels for our analysis.
  \item We can use the \code{levels()} function to check the levels:
\end{itemize}
\begin{lstlisting}[language=R]
levels(f1)
\end{lstlisting}
\begin{lstlisting}[style=output]
[1] "0" "1"
\end{lstlisting}
\begin{itemize}
  \item \code{as.factor()} will have the same result as \code{factor()}. This is preferred because it will provide a quicker return.
\end{itemize}
\begin{lstlisting}[language=R]
f2 <- as.factor(c("Female", "Male", "Female", "Unknown", "Female", "Male", "Male", "Female"))
f2
\end{lstlisting}
\begin{lstlisting}[style=output]
[1] Female  Male    Female  Unknown Female  Male    Male    Female 
Levels: Female Male Unknown
\end{lstlisting}
\end{frame}

\begin{frame}[fragile]{Tables}
\begin{lstlisting}[language=R]
f2
\end{lstlisting}
\begin{lstlisting}[style=output]
[1] Female  Male    Female  Unknown Female  Male    Male    Female 
Levels: Female Male Unknown
\end{lstlisting}
\begin{itemize}
  \item \code{table()} will produce a table to show how often each object for a categorical (factor) variable is given:
\end{itemize}
\begin{lstlisting}[language=R]
table(f2)
\end{lstlisting}
\begin{lstlisting}[style=output]
f2
 Female    Male Unknown 
      4       3       1
\end{lstlisting}
\begin{itemize}
  \item Small trick: to show relative frequency \textit{(or percentage)}, we can do the following:
\end{itemize}
\begin{lstlisting}[language=R]
table(f2) / length(f2)
\end{lstlisting}
\begin{lstlisting}[style=output]
f2
 Female    Male Unknown 
  0.500   0.375   0.125
\end{lstlisting}
\end{frame}

\begin{frame}[fragile]{Tables (cont.)}
\begin{lstlisting}[language=R]
gender <- table(f2) / length(f2) # the result will be a named vector
\end{lstlisting}
\begin{columns}[T]
\begin{column}{0.48\textwidth}
\begin{lstlisting}[language=R]
gender[1]
\end{lstlisting}
\begin{lstlisting}[style=output]
Female 
   0.5
\end{lstlisting}
\end{column}
\begin{column}{0.48\textwidth}
\begin{lstlisting}[language=R]
names(gender)
\end{lstlisting}
\begin{lstlisting}[style=output]
[1] "Female"  "Male"    "Unknown"
\end{lstlisting}
\end{column}
\end{columns}
\begin{itemize}
  \item Two-way table using \code{table()}
\end{itemize}
\begin{lstlisting}[language=R]
exam1 <- as.factor(c("A", "B", "F", "D", "B", "C", "A"))
exam2 <- as.factor(c("A", "B", "C", "A", "C", "B", "B"))

table(exam1, exam2) # a two-way table
\end{lstlisting}
\begin{lstlisting}[style=output]
     exam2
exam1 A B C
    A 1 1 0
    B 0 1 1
    C 0 1 0
    D 1 0 0
    F 0 0 1
\end{lstlisting}
\end{frame}

\begin{frame}[fragile]{File Input/Output}
\begin{itemize}
  \item \textit{Save the four reference files to your class folder.} \textit{NOTE: for the .txt file, you need to right-click and then select “save link as”.}
  \item There are multiple ways to import \& export data, depending on what type of file/database it is and where it is stored.
  \item \textbf{Connecting to text files:} \code{read.table()} is used to import data from \code{.txt} files:
\end{itemize}
\begin{lstlisting}[language=R]
grade_df <- read.table("R Ref Files/grade.txt", header = TRUE)
\end{lstlisting}
\end{frame}

\begin{frame}[fragile]{File Input/Output (cont.)}
\begin{columns}[T]
\begin{column}{0.48\textwidth}
\begin{lstlisting}[language=R]
class(grade_df$Name)
\end{lstlisting}
\begin{lstlisting}[style=output]
[1] "character"
\end{lstlisting}
\end{column}
\begin{column}{0.48\textwidth}
\begin{lstlisting}[language=R]
class(grade_df$Grade)
\end{lstlisting}
\begin{lstlisting}[style=output]
[1] "integer"
\end{lstlisting}
\end{column}
\end{columns}
\begin{itemize}
  \item Using the \code{stringsAsFactors} argument \textit{(by default = FALSE)}:
\end{itemize}
\begin{lstlisting}[language=R]
grade_df2 <- read.table("R Ref Files/grade.txt", header = TRUE, stringsAsFactors = TRUE)
\end{lstlisting}
\begin{lstlisting}[language=R]
class(grade_df2$Name)
\end{lstlisting}
\begin{lstlisting}[style=output]
[1] "factor"
\end{lstlisting}
\end{frame}

\begin{frame}[fragile]{File Input/Output (cont.)}
\begin{itemize}
  \item \textbf{Connecting to CSVs:} \code{read.csv()} is used to import data from \code{.csv} file:
\end{itemize}
\begin{lstlisting}[language=R]
tl_df <- read.csv("R Ref Files/termlife.csv", header = TRUE)

# convert the categorical variables in tl_df into factor variables
tl_df$gender <- as.factor(tl_df$gender)
tl_df$marstat <- as.factor(tl_df$marstat)
tl_df$edu <- as.factor(tl_df$edu)
tl_df$numhh <- as.factor(tl_df$numhh)
\end{lstlisting}
\begin{itemize}
  \item \textbf{Connecting to Excel:} first, you need to install the \code{readxl} package. \textit{We’ll talk more about R packages in our next class.}
  \item Then, you can use functions from the package to import your excel data.
\end{itemize}
\end{frame}

\begin{frame}[fragile]{File Input/Output (cont.)}
\begin{lstlisting}[language=R]
# Import Excel data and name your data in R as follows:
library(readxl)
# To work properly with the below function, the Excel table
# needs to start in cell A1 of the respective worksheet
excel_data <- read_excel(path="C:/Users/ssaul/Desktop/MEPS_year.xlsx", 
                         sheet = "raw_data")
\end{lstlisting}
\begin{itemize}
  \item You can also import SAS files and connect to other servers. For example, in many companies, it will be helpful to connect directly to Oracle servers.
  \item \textbf{Exporting data:} \code{write.table()} and \code{write.csv()} can be used to export \code{data.frame} to files:
\end{itemize}
\begin{lstlisting}[language=R]
write.csv(grade_df, "R Ref Files/grade.csv")
\end{lstlisting}
\end{frame}

